برای حذف نویز تناوبی موجود در تصویر، طراحی فیلتر در حوزه فرکانس به صورت \lr{notch-reject} مناسب است؛ چون در طیف، نویزهای تناوبی به شکل قله‌های نقطه‌ای متقارن ظاهر می‌شوند و تصویر اصلی عمدتا در مرکز متمرکز است.

مراحل کار:
\begin{flalign*}
&\text{۱) تبدیل فوریه و مرکز کردن:}\quad F(u,v)=\mathcal{F}\{f(x,y)\},\;\;F_c=\text{fftshift}(F). \\
&\text{۲) مکان‌یابی قله‌های نویز:}\quad (u_k,v_k)\;\; \text{و قرینه‌ها}\;(-u_k,-v_k). \\
&\text{۳) ساخت فیلتر \lr{notch-reject} گاوسی:} \\
&H(u,v)=\prod_{k=1}^{K}\left(1-\exp\left(-\frac{D_k^2(u,v)}{2\sigma^2}\right)\right)
\left(1-\exp\left(-\frac{D_k'^2(u,v)}{2\sigma^2}\right)\right), \\
&\quad D_k^2(u,v)=(u-u_k)^2+(v-v_k)^2,\;\;D_k'^2(u,v)=(u+u_k)^2+(v+v_k)^2. \\
&\text{۴) اعمال فیلتر و بازگشت به حوزه مکان:} \\
&G_c(u,v)=H(u,v)\,F_c(u,v),\quad g(x,y)=\Re\{\mathcal{F}^{-1}(\text{ifftshift}(G_c))\}. \\
\end{flalign*}

توضیح انتخاب:
\begin{flalign*}
&\text{قله‌های روشنِ متقارن در طیف نشان‌دهنده نویز تناوبی هستند؛}\\
&\text{فیلتر \lr{notch-reject} فقط همان فرکانس‌ها را حذف می‌کند و جزئیات تصویر حفظ می‌شود.} \\
&\text{پارامتر } \sigma \text{ با پهنای قله‌ها تنظیم می‌شود؛ کوچک‌تر برای حذف دقیق‌تر.}
\end{flalign*}

در صورت وجود خطوط روشن در طیف (نویز نواری)، می‌توان علاوه بر نقاط، یک \lr{notch band-reject} باریک روی همان خطوط نیز افزود تا شبکه‌های سراسری حذف شوند.